% Part: normal-modal-logic
% Chapter: completeness
% Section: complete-consistent-sets

\documentclass[../../../include/open-logic-section]{subfiles}

\begin{document}

\olfileid[zh]{nml}{com}{ccs}

\olsection{完全 $\Sigma$-一致集}

设 $\Sigma$ 是一组模态!!{formula}的集合——可将它们看作某一正规模态逻辑的公理或定义原则。一个集合 $\Gamma$ 是 $\Sigma$-一致的，当且仅当 $\Gamma \Proves/[\Sigma] \lfalse$，即不存在从 $\Sigma$ 出发、形如 $!A_1 \lif (!A_2 \lif \cdots (!A_n \lif \lfalse)\dots)$ 的!!{derivation}，其中每个 $!A_i \in \Gamma$。我们将构造一个「典范」模型，其中每个世界都被视为一种特殊的 $\Sigma$-一致集：它不只是 $\Sigma$-一致的，而是极大一致的，即在它决定每个模态!!{formula}的真值这一意义上：对每个 $!A$，要么 $!A \in \Gamma$，要么 $\lnot !A \in \Gamma$：

\begin{defn}
  一个集合 $\Gamma$ 是\emph{完全 $\Sigma$-一致的（complete $\Sigma$-consistent）}，当且仅当它是 $\Sigma$-一致的，且对每个~$!A$，要么 $!A \in \Gamma$，要么 $\lnot !A \in \Gamma$。
\end{defn}

完全 $\Sigma$-一致集 $\Gamma$ 具有若干有用的性质。首先，它们是演绎封闭的，即若 $\Gamma \Proves[\Sigma] !A$ 则 $!A \in \Gamma$。这尤其意味着，!!a{formula}~$!A \in \Sigma$ 的每个特例亦 $\in \Gamma$。再者，$\Gamma$ 中的成员关系镜像了命题联结词的真值条件。这在我们定义「典范模型」时将十分重要。

\begin{prop}\ollabel{prop:ccs-properties}
  设 $\Gamma$ 是完全 $\Sigma$-一致的。则：
  \begin{enumerate}
  \item \ollabel{prop:ccs-closed}%
    $\Gamma$ 在~$\Sigma$ 中是演绎封闭的（deductively closed）。
  \item \ollabel{prop:ccs-sigma}%
    $\Sigma \subseteq \Gamma$。
  \tagitem{prvFalse}{\ollabel{prop:ccs-lfalse}%
    $\lfalse \notin \Gamma$}{}
  \tagitem{prvTrue}{\ollabel{prop:ccs-ltrue}%
    $\ltrue \in \Gamma$}{}
  \item \ollabel{prop:ccs-lnot}%
    $\lnot!A \in \Gamma$ 当且仅当 $!A \notin \Gamma$。
  \tagitem{prvAnd}{\ollabel{prop:ccs-land}%
    $!A \land !B \in \Gamma$ 当且仅当 $!A \in \Gamma$ 且 $!B \in \Gamma$}{}
  \tagitem{prvOr}{\ollabel{prop:ccs-lor}%
    $!A \lor !B \in \Gamma$ 当且仅当 $!A \in \Gamma$ 或 $!B \in \Gamma$}{}
  \tagitem{prvIf}{\ollabel{prop:ccs-lif}%
    $!A \lif !B \in \Gamma$ 当且仅当 $!A \notin \Gamma$ 或 $!B \in \Gamma$}{}
  \tagitem{prvIff}{\ollabel{prop:ccs-liff}%
    $!A \liff !B \in \Gamma$ 当且仅当要么 $!A \in \Gamma$ 且 $!B \in \Gamma$，要么 $!A \notin \Gamma$ 且 $!B \notin \Gamma$}{}
  \end{enumerate}
\end{prop}

\begin{proof}
  \begin{enumerate}
  \item 设 $\Gamma \Proves[\Sigma] !A$ 但 $!A \notin \Gamma$。则由于 $\Gamma$ 是完全 $\Sigma$-一致的，故 $\lnot!A \in \Gamma$。这将使 $\Gamma$ 不一致，因为 $!A, \lnot !A \Proves[\Sigma] \lfalse$。

  \item 若 $!A \in \Sigma$ 则 $\Gamma \Proves[\Sigma] !A$，借助演绎封闭性（即情形~\olref{prop:ccs-closed}），有 $!A \in \Gamma$。

  \tagitem{prvFalse}{若 $\lfalse \in \Gamma$，则 $\Gamma \Proves[\Sigma] \lfalse$，于是 $\Gamma$ 将是 $\Sigma$-不一致的。}{}

  \tagitem{prvTrue}{$\Gamma \Proves[\Sigma] \ltrue$，故借助演绎封闭性（即情形~\olref{prop:ccs-closed}），有 $\ltrue \in \Gamma$。}{}

\item 若 $\lnot!A \in \Gamma$，则据一致性有 $!A \notin \Gamma$；而若 $!A \notin \Gamma$，则因 $\Gamma$ 是完全 $\Sigma$-一致的，故 $!A \in \Gamma$。

  \tagitem{prvAnd}{\iftag{probAnd}{习题。}{设 $!A \land !B \in \Gamma$。由于 $(!A \land !B) \lif !A$ 是重言式特例，故借助演绎封闭性（即情形~\olref{prop:ccs-closed}），有 $!A \in \Gamma$。对 $!B \in \Gamma$ 同理。另一方面，设 $!A \in \Gamma$ 且 $!B \in \Gamma$。则演绎封闭性蕴含 $(!A \land !B) \in \Gamma$，因为 $!A \lif (!B \lif (!A \land !B))$ 是重言式特例。}}{}

  \tagitem{prvOr}{\iftag{probOr}{习题。}{设 $!A \lor !B \in \Gamma$，且 $!A \notin \Gamma$ 且 $!B \notin \Gamma$。由于 $\Gamma$ 是完全 $\Sigma$-一致的，故 $\lnot!A \in \Gamma$ 且 $\lnot !B \in \Gamma$。则 $\lnot(!A \lor !B) \in \Gamma$，因为 $\lnot !A \lif (\lnot !B \lif \lnot (!A \lor !B))$ 是重言式特例。这将意味着 $\Gamma$ 是 $\Sigma$-不一致的，即矛盾。}}{}

  \tagitem{prvIf}{\iftag{probIf}{习题。}{设 $!A \lif !B \in \Gamma$ 且 $!A \in \Gamma$；则 $\Gamma \Proves[\Sigma] !B$，从而借助演绎封闭性有 $!B \in \Gamma$。反之，若 $!A \lif !B \notin \Gamma$，则因 $\Gamma$ 是完全 $\Sigma$-一致的，故 $\lnot (!A \lif !B) \in \Gamma$。由于 $\lnot(!A \lif !B) \lif !A$ 是重言式特例，故借助演绎封闭性有 $!A \in \Gamma$。由于 $\lnot(!A \lif !B) \lif \lnot !B$ 是重言式特例，故 $\lnot !B \in \Gamma$。则因 $\Gamma$ 是 $\Sigma$-一致的，故 $!B \notin \Gamma$。}}{}

  \tagitem{prvIff}{\iftag{probIff}{习题。}{设 $!A \liff !B \in \Gamma$。若 $!A \in \Gamma$，则 $!B \in \Gamma$，因为 $(!A \liff !B) \lif (!A \lif !B)$ 是重言式特例。类似地，若 $!B \in \Gamma$，则 $!A \in \Gamma$。于是要么 $!A \in \Gamma$ 且 $!B \in \Gamma$，要么 $!A \notin \Gamma$ 且 $!B \notin \Gamma$。

      反之，设 $!A \lif !B \notin \Gamma$。由于 $\Gamma$ 是完全 $\Sigma$-一致的，故 $\lnot (!A \liff !B) \in \Gamma$。由于 $\lnot(!A \liff !B) \lif (!A \lif \lnot !B)$ 是重言式特例，若 $!A \in \Gamma$ 则 $\lnot !B \in \Gamma$，且因 $\Gamma$ 是 $\Sigma$-一致的，故 $!B \notin \Gamma$。类似地，若 $!B \in \Gamma$ 则 $!A \notin \Gamma$。于是既非 $!A \in \Gamma$ 且 $!B \in \Gamma$，也非 $!A \notin \Gamma$ 且 $!B \notin \Gamma$。}}{}
  \end{enumerate}
\end{proof}

\begin{probtag}{probAnd,probOr,probIf,probIff}
  完成 \olref[nml][com][ccs]{prop:ccs-properties} 的证明。
\end{probtag}

\end{document}
